\section{Approximate original gate expansions}
\label{sec:peps_approximate_original_gates}

The final paragraph of the proof of
\cite[Theorem~5.2]{OpenAI2026PolynomialPEPS} permits approximate
monomial expansions. We retain the original gate occurrences and rescale
only their scalar coefficients. The resulting operator and physical
trace-norm estimates do not involve private dimensions.

\begin{definition}[Original operators on a fixed chronology]
    \label{def:peps_approximate_original_operators}
    \lean{TNLean.PEPS.PairEffect.EffectCircuit.GateMaps}
    \lean{TNLean.PEPS.PairEffect.EffectCircuit.evalWithGateMaps}
    \lean{TNLean.PEPS.PairEffect.EffectCircuit.IsGateApproximation}
    \leanok
    Fix a chronological circuit whose gate occurrence $g$ carries an
    ordered expansion $T_g=\sum_j c_{g,j}M_{g,j}$. Assign an original
    operator $G_g$ on the same local input and output memories to every
    occurrence. Let $W_G$ be their chronological composition, with the
    prescribed private operations and spectator identities.
    For $\delta\ge0$, require all private operators and all $G_g$ to be
    contractions, every listed monomial to be allowed, and
    $\|T_g-G_g\|\le\delta$. Retain the convention that a nonprivate
    occurrence has a participating set of cardinality different from one;
    the empty participating set is permitted.
\end{definition}

\begin{theorem}[Contractivity of the original chronology]
    \label{thm:peps_original_operator_contraction}
    \lean{TNLean.PEPS.PairEffect.EffectCircuit.norm_evalWithGateMaps_le_one}
    \lean{TNLean.PEPS.PairEffect.OriginalCircuit.norm_eval_le_one}
    \leanok
    \uses{def:peps_approximate_original_operators}
    Under these local hypotheses, $\|W_G\|\le1$. The same assertion
    holds for every original contraction circuit.
\end{theorem}
\begin{proof}
    \leanok
    Induct on the chronology. Composition multiplies operator norms,
    and tensoring with a spectator identity does not increase the norm.
    The prescribed owner and register identifications are isometries.
\end{proof}

\begin{definition}[Rescaling an approximate expansion]
    \label{def:peps_rescaled_original}
    \lean{TNLean.PEPS.PairEffect.rescalePartyGate}
    \lean{TNLean.PEPS.PairEffect.EffectCircuit.rescaledOriginal}
    \leanok
    \uses{def:peps_approximate_original_operators}
    Replace each coefficient $c_{g,j}$ by $c_{g,j}/(1+\delta)$.
    Write $\widetilde W_\delta$ for the chronological circuit with these
    new expansions and the original private operations. Thus its local
    operator at $g$ is $\widetilde T_g=T_g/(1+\delta)$.
\end{definition}
\begin{proof}
    \leanok
    \lean{TNLean.PEPS.PairEffect.monomials_rescalePartyGate}
    \lean{TNLean.PEPS.PairEffect.gate_rescalePartyGate}
    \lean{TNLean.PEPS.PairEffect.sum_norm_rescalePartyGate_le}
    \lean{TNLean.PEPS.PairEffect.isAllowed_rescalePartyGate}
    The ordered monomial list is unchanged. Linearity gives the displayed
    local operator, and $\|T_g\|\le\|G_g\|+\delta\le1+\delta$
    proves its contractivity. Each monomial remains allowed, while
    \begin{align}
        \sum_j\left|\frac{c_{g,j}}{1+\delta}\right|
        \le\sum_j|c_{g,j}|.
        \label{eq:peps_rescaled_coefficient_sum}
    \end{align}
\end{proof}

\begin{definition}[Occurrences and monomial counts after rescaling]
    \label{def:peps_rescaled_occurrences}
    \lean{TNLean.PEPS.PairEffect.EffectCircuit.rescaledLocationsEquiv}
    \lean{TNLean.PEPS.PairEffect.EffectCircuit.monomialCount}
    \leanok
    \uses{def:peps_rescaled_original}
    Match every gate of $\widetilde W_\delta$ to the same original
    occurrence, and let $k_g$ be the length of its prescribed monomial
    list. Repeated occurrences remain distinct in this correspondence.
\end{definition}

\begin{theorem}[Preservation of the original resource bounds]
    \label{thm:peps_rescaled_resources}
    \lean{TNLean.PEPS.PairEffect.EffectCircuit.nonprivateCount_rescaledOriginal}
    \lean{TNLean.PEPS.PairEffect.EffectCircuit.participants_rescaledOriginal}
    \lean{TNLean.PEPS.PairEffect.EffectCircuit.participationCount_rescaledOriginal}
    \lean{TNLean.PEPS.PairEffect.EffectCircuit.isExpansionBounded_rescaledOriginal}
    \lean{TNLean.PEPS.PairEffect.EffectCircuit.isMonomialBounded_rescaledOriginal}
    \leanok
    \uses{def:peps_rescaled_occurrences}
    Rescaling preserves the number $N$ of nonprivate occurrences, their
    participating sets, and each party's number of participations.
    If each monomial has at most $r$ pair effects, each gate has
    absolute coefficient sum at most $S$, and $k_g\le K$, the same
    three bounds hold for $\widetilde W_\delta$.
\end{theorem}
\begin{proof}
    \leanok
    \uses{def:peps_rescaled_original,def:peps_rescaled_occurrences}
    The occurrence correspondence preserves each participating set.
    The monomials and their pair effects are identical, and
    (\ref{eq:peps_rescaled_coefficient_sum}) bounds the new coefficient sums.
\end{proof}

\begin{theorem}[Two-factor telescoping estimate]
    \label{thm:peps_contraction_composition_error}
    \lean{TNLean.PEPS.PairEffect.norm_comp_sub_comp_le_one}
    \leanok
    Let $A,C:E\to F$ and $B,D:F\to G$ be bounded linear maps between
    complex inner-product spaces. If $\|B\|,\|C\|\le1$, then
    $\|BA-DC\|\le\|A-C\|+\|B-D\|$.
\end{theorem}
\begin{proof}
    \leanok
    Use $BA-DC=B(A-C)+(B-D)C$, the triangle inequality, and
    submultiplicativity of the operator norm.
\end{proof}

\begin{theorem}[Chronological operator error]
    \label{thm:peps_approximate_chronological_error}
    \lean{TNLean.PEPS.PairEffect.EffectCircuit.norm_rescaledOriginal_sub_evalWithGateMaps_le}
    \leanok
    \uses{def:peps_rescaled_original}
    If the local approximation error is at most $\delta\ge0$ and the
    original chronology has $N$ nonprivate occurrences, then
    $\|\widetilde W_\delta-W_G\|\le2\delta N$.
\end{theorem}
\begin{proof}
    \leanok
    \uses{thm:peps_original_operator_contraction,thm:peps_contraction_composition_error}
    The local error satisfies
    \begin{align}
        \|\widetilde T_g-G_g\|
        \le \frac{\|T_g-G_g\|+\delta\|G_g\|}{1+\delta}
        \le2\delta.
        \label{eq:peps_rescaled_local_error}
    \end{align}
    Apply the two-factor estimate at every composition. Private maps and
    register permutations add no error; spectator identities preserve
    the local bound. Induction adds one contribution for each gate.
\end{proof}

\begin{definition}[Physical output of the original operators]
    \label{def:peps_original_gate_density}
    \lean{TNLean.PEPS.PairEffect.EffectCircuit.physicalDensityWithGateMaps}
    \leanok
    \uses{def:peps_approximate_original_operators}
    Choose an isometric identification $K$ of the final memory with
    finite physical and discarded coordinate spaces. For an input $x$, set
    $\rho_G(x)=\operatorname{Tr}_{\rm discard}|KW_Gx\rangle\langle KW_Gx|$.
    Define $\widetilde\rho_\delta(x)$ by the same formula with
    $\widetilde W_\delta$. These operators are not normalized by their trace.
\end{definition}

\begin{theorem}[Physical error of approximate original gates]
    \label{thm:peps_approximate_physical_error}
    \lean{TNLean.PEPS.PairEffect.EffectCircuit.rectangularTraceNorm_rescaledOriginal_density_sub_le}
    \lean{TNLean.PEPS.PairEffect.EffectCircuit.rectangularTraceNorm_rescaledOriginal_density_sub_le_budget}
    \leanok
    \uses{def:peps_original_gate_density}
    If $\|x\|\le1$, then
    $\|\widetilde\rho_\delta(x)-\rho_G(x)\|_1\le4\delta N$.
    In particular, for $\varepsilon>0$ and $N\le M$, taking
    $\delta=\varepsilon/(8\max\{1,M\})$ gives error at most
    $\varepsilon/2$. Replacing this budget parameter by
    $\varepsilon/2$ reserves $\varepsilon/4$ for this step.
\end{theorem}
\begin{proof}
    \leanok
    \uses{thm:peps_approximate_chronological_error,thm:peps_original_operator_contraction}
    Put $u=K\widetilde W_\delta x$ and $v=KW_Gx$.
    Contractivity and the isometry $K$ give $\|u\|,\|v\|\le1$.
    Therefore
    \begin{align}
        \|\widetilde\rho_\delta(x)-\rho_G(x)\|_1
         & =\bigl\|\operatorname{Tr}_{\rm discard}
        (|u\rangle\langle u|-|v\rangle\langle v|)\bigr\|_1
        \le2\|u-v\|                                \\
         & \le2\|\widetilde W_\delta-W_G\|\,\|x\|
        \le4\delta N.
        \label{eq:peps_approximate_density_error}
    \end{align}
    Substitute the prescribed scalar budget into
    (\ref{eq:peps_approximate_density_error}).
\end{proof}

\begin{theorem}[Natural inverse-power accuracy]
    \label{thm:peps_approximate_natural_accuracy}
    \lean{TNLean.PEPS.PairEffect.gateApproximation_error_inv_pow_le}
    \lean{TNLean.PEPS.PairEffect.EffectCircuit.rectangularTraceNorm_rescaledOriginal_density_sub_le_inv_pow}
    \uses{def:peps_original_gate_density}
    \leanok
    Let $N\le CL^n$, where $n,p$ are nonnegative integers and
    $L\ge\max\{1,16C\}$. Local accuracy
    $\delta=L^{-(n+p+1)}$ gives
    $\|\widetilde\rho_\delta(x)-\rho_G(x)\|_1\le L^{-p}/4$
    for $\|x\|\le1$.
\end{theorem}
\begin{proof}
    \leanok
    \uses{thm:peps_approximate_physical_error}
    The scalar estimate is
    $4L^{-(n+p+1)}N\le(16C/L)L^{-p}/4\le L^{-p}/4$.
\end{proof}

\begin{theorem}[One fixed accuracy exponent for approximate gates]
    \label{thm:peps_approximate_gate_real_accuracy}
    \lean{TNLean.PEPS.PairEffect.gateApproximation_error_real_pow_le}
    \leanok
    Let $L\ge2$, let $N,n$ be nonnegative integers, and suppose
    $N\le C L^n$. For a real target exponent $p$, set
    \begin{align}
        q=n+\lceil p\rceil_++\lceil16C\rceil_+,
        \qquad \lceil x\rceil_+=\max(0,\lceil x\rceil).
    \end{align}
    Then
    \begin{align}
        4L^{-q}N\le \frac14L^{-p}.
    \end{align}
    In particular, $q$ depends on the fixed polynomial bound and on $p$,
    but not on $L$.
\end{theorem}
\begin{proof}
    \leanok
    Write $s=\lceil16C\rceil_+$. The elementary inequalities
    $16C\le s\le2^s\le L^s$ imply
    \begin{align}
        4L^{-q}N\le \frac{16C}{L^s}\,
        \frac{L^{-\lceil p\rceil_+}}4
        \le \frac{L^{-p}}4.
    \end{align}
    The argument includes $N=0$ and requires no further size threshold.
\end{proof}

\begin{theorem}[Approximate original gates at real inverse-power accuracy]
    \label{thm:peps_approximate_original_real_accuracy}
    \lean{TNLean.PEPS.PairEffect.EffectCircuit.rectangularTraceNorm_rescaledOriginal_density_sub_le_real_pow}
    \leanok
    Suppose a chronological contraction circuit has $N\le C L^n$ nonprivate
    gate occurrences, where $L\ge2$. Fix a real exponent $p$, and choose $q$
    as in Theorem~\ref{thm:peps_approximate_gate_real_accuracy}.
    At every original gate, let an allowed monomial expansion approximate
    that gate in operator norm to accuracy $L^{-q}$. Divide all coefficients
    of this expansion by $1+L^{-q}$. For an input vector of norm at most one,
    the physical output of the resulting contraction circuit differs from
    the original physical density by at most $L^{-p}/4$ in trace norm.
    The output density is not normalized; all private registers are traced out.
\end{theorem}
\begin{proof}
    \leanok
    \uses{thm:peps_approximate_gate_real_accuracy,thm:peps_approximate_physical_error}
    The chronological physical-density estimate is $4L^{-q}N$.
    Apply Theorem~\ref{thm:peps_approximate_gate_real_accuracy}. The comparison uses the actual original gate
    operators, the same private operations and the same physical readout.
\end{proof}

\begin{definition}[Termwise bounds on the original gate expansions]
    \label{def:peps_original_termwise_bounds}
    \lean{TNLean.PEPS.PairEffect.OriginalCircuit.IsTermwiseBounded}
    \leanok
    An original contraction circuit has termwise bounds $(r,C)$ if, at
    every nonprivate gate occurrence, each monomial in its prescribed
    expansion contains at most $r$ pair effects and has scalar coefficient
    of modulus at most $C$. These conditions concern the actual ordered
    expansions. They do not impose a bound on the sum of the moduli.
    The constituent local maps already satisfy the contraction conditions
    of the original circuit.
\end{definition}

\begin{theorem}[Coefficient sums of the actual gate lists]
    \label{thm:peps_gate_coefficient_sum}
    \lean{TNLean.PEPS.PairEffect.PartyGate.sum_norm_coefficients_le}
    \lean{TNLean.PEPS.PairEffect.PartyGate.sum_norm_coefficients_le_of_power_bounds}
    \leanok
    Write a prescribed gate expansion as $\sum_{j=1}^{k}c_jM_j$.
    If $|c_j|\leq C$ for every listed term, then
    $\sum_{j=1}^{k}|c_j|\leq kC$.
    In particular, if $L\geq0$, $C_C\geq0$, $u,v$ are nonnegative integers,
    $k\leq C_KL^u$, and $|c_j|\leq C_CL^v$ for every listed term, then
    \begin{align}
        \sum_{j=1}^{k}|c_j|\leq C_KC_CL^{u+v}.
        \label{eq:peps_termwise_polynomial_sum}
    \end{align}
    Both statements include the empty expansion. This is the coefficient
    estimate used in Theorem~5.2, assumption~3, and the opening reduction
    in its proof.
\end{theorem}
\begin{proof}
    \leanok
    Sum the individual inequalities over the ordered list. For the
    polynomial estimate, multiply its length bound by the nonnegative
    number $C_CL^v$ and combine the powers of $L$.
\end{proof}

\begin{theorem}[Uniform coefficient sums for the original circuit]
    \label{thm:peps_original_coefficient_sum}
    \lean{TNLean.PEPS.PairEffect.OriginalCircuit.isExpansionBounded_of_isTermwiseBounded}
    \lean{TNLean.PEPS.PairEffect.OriginalCircuit.isExpansionBounded_of_termwise_power_bounds}
    \uses{def:peps_original_termwise_bounds}
    \leanok
    Suppose an original circuit has at most $K$ monomials at each
    nonprivate gate occurrence and has termwise bounds $(r,C)$, where
    $C\geq0$. Then every such gate has absolute coefficient sum at most
    $KC$, or at most any $S\geq KC$, and each monomial still has at most
    $r$ pair effects.
    Consequently, if $K\leq C_KL^u$ and the termwise bounds are
    $(r,C_CL^v)$, with $L,C_C\geq0$, the uniform coefficient-sum bound is
    $C_KC_CL^{u+v}$. No nonemptiness of the expansion or participating
    set is required.
\end{theorem}
\begin{proof}
    \leanok
    \uses{def:peps_original_termwise_bounds,thm:peps_gate_coefficient_sum}
    Apply the preceding list estimate at each original gate occurrence.
    Composition and addition of spectator registers retain these same
    local conditions; private operations and register permutations add
    no nonprivate gate expansions. The effect count is unchanged.
\end{proof}
